%% ============================================================
%% app_ohm.tex — Appendice OpenHistoryMap
%% ============================================================

\chapter{Appendice OpenHistoryMap}
\markboth{OpenHistoryMap}{OpenHistoryMap}

\lettrine[lines=3]{Q}{uesta} appendice fornisce i dati geografici
storici di \textit{Vespri 1282} in formato RDF/Turtle compatibile
con \textbf{OpenHistoryMap} (\url{https://openhistorymap.org}).
I dati coprono Palermo in dettaglio e i principali centri siciliani
in forma ridotta, e sono progettati come risorsa standalone
espandibile indipendentemente dal gioco.

I dati usano il sistema di coordinate WGS84. Le posizioni sono
approssimazioni storiche basate sulla storiografia disponibile
e sull'analisi delle planimetrie medievali.

\section{Prefissi}

\begin{verbatim}
@prefix rdf:    <http://www.w3.org/1999/02/22-rdf-syntax-ns#> .
@prefix rdfs:   <http://www.w3.org/2000/01/rdf-schema#> .
@prefix owl:    <http://www.w3.org/2002/07/owl#> .
@prefix xsd:    <http://www.w3.org/2001/XMLSchema#> .
@prefix geo:    <http://www.w3.org/2003/01/geo/wgs84_pos#> .
@prefix geosparql: <http://www.opengis.net/ont/geosparql#> .
@prefix ohm:    <https://openhistorymap.org/ontology/> .
@prefix ohmplace: <https://openhistorymap.org/place/> .
@prefix v1282:  <https://openhistorymap.org/vespri1282/> .
@prefix time:   <http://www.w3.org/2006/time#> .
@prefix dcterms: <http://purl.org/dc/terms/> .
\end{verbatim}

\section{Periodo Storico}

\begin{verbatim}
v1282:period_1282 a ohm:HistoricalPeriod ;
  rdfs:label "Palermo 1282 — Vespri Siciliani"@it ,
             "Palermo 1282 — Sicilian Vespers"@en ;
  time:hasBeginning [ time:inXSDDate "1282-02-01"^^xsd:date ] ;
  time:hasEnd       [ time:inXSDDate "1282-03-31"^^xsd:date ] ;
  ohm:politicalContext "Dominio angioino / Angevin rule"@it .
\end{verbatim}

\section{Palermo — Luoghi Principali}

\begin{verbatim}

%% --- QUARTIERI ---

ohmplace:palermo_cassaro a ohm:HistoricalQuarter ;
  rdfs:label "Cassaro"@it , "Cassaro"@en ;
  dcterms:description
    "Centro storico di Palermo, sede del Palazzo dei Normanni.
     Deriva dall'arabo al-Qasr (il castello)."@it ,
    "Historic centre of Palermo, site of the Norman Palace.
     Name derives from Arabic al-Qasr (the castle)."@en ;
  geo:lat "38.1115"^^xsd:decimal ;
  geo:long "13.3523"^^xsd:decimal ;
  ohm:period v1282:period_1282 ;
  ohm:ethnicity "Normanna/Francese"@it , "Norman/French"@en ;
  v1282:narrativeRole "Sede del potere angioino"@it .

ohmplace:palermo_kalsa a ohm:HistoricalQuarter ;
  rdfs:label "Kalsa"@it , "Kalsa"@en ;
  dcterms:description
    "Quartiere arabo di Palermo, vicino al porto.
     Dall'arabo al-Khalisa (la pura, l'eletta)."@it ,
    "Arab quarter of Palermo, near the port.
     From Arabic al-Khalisa (the pure, the chosen)."@en ;
  geo:lat "38.1132"^^xsd:decimal ;
  geo:long "13.3671"^^xsd:decimal ;
  ohm:period v1282:period_1282 ;
  ohm:ethnicity "Araba/Mista"@it , "Arab/Mixed"@en ;
  v1282:narrativeRole "Porto e commercio"@it .

ohmplace:palermo_seralcadio a ohm:HistoricalQuarter ;
  rdfs:label "Seralcadio"@it , "Seralcadio"@en ;
  dcterms:description
    "Quartiere degli artigiani e mercanti."@it ,
    "Quarter of craftsmen and merchants."@en ;
  geo:lat "38.1145"^^xsd:decimal ;
  geo:long "13.3544"^^xsd:decimal ;
  ohm:period v1282:period_1282 ;
  v1282:narrativeRole "Artigiani — fazione Popolo"@it .

ohmplace:palermo_albergheria a ohm:HistoricalQuarter ;
  rdfs:label "Albergheria"@it , "Albergheria"@en ;
  dcterms:description
    "Il quartiere più povero e densamente popolato di Palermo."@it ,
    "The poorest and most densely populated quarter of Palermo."@en ;
  geo:lat "38.1118"^^xsd:decimal ;
  geo:long "13.3545"^^xsd:decimal ;
  ohm:period v1282:period_1282 ;
  v1282:narrativeRole "Ceto popolare — fazione Popolo"@it .

%% --- EDIFICI E LUOGHI CHIAVE ---

ohmplace:palermo_santo_spirito a ohm:HistoricalBuilding ;
  rdfs:label "Chiesa di Santo Spirito"@it ,
             "Church of Santo Spirito"@en ;
  dcterms:description
    "Chiesa normanna fuori le mura di Palermo. Teatro della
     processione del 30 marzo 1282 e scintilla della rivolta.
     Nota anche come 'chiesa dei Vespri'."@it ,
    "Norman church outside the walls of Palermo. Site of the
     Easter Monday procession of 30 March 1282 and spark of
     the uprising. Also known as the 'church of the Vespers'."@en ;
  geo:lat "38.1048"^^xsd:decimal ;
  geo:long "13.3595"^^xsd:decimal ;
  ohm:period v1282:period_1282 ;
  ohm:buildingType "chiesa"@it , "church"@en ;
  ohm:constructionPeriod "XII sec."@it , "12th century"@en ;
  v1282:narrativeRole "Detonatore della rivolta"@it ,
                       "Trigger of the uprising"@en ;
  v1282:faction v1282:faction_clero ;
  owl:sameAs <https://www.wikidata.org/entity/Q3960295> .

ohmplace:palermo_normanni a ohm:HistoricalBuilding ;
  rdfs:label "Palazzo dei Normanni"@it ,
             "Norman Palace"@en ;
  dcterms:description
    "Sede del potere angioino a Palermo. Costruito in origine
     dagli emiri arabi, ampliato dai Normanni, usato dagli
     Svevi e poi dagli Angioini."@it ,
    "Seat of Angevin power in Palermo. Originally built by Arab
     emirs, enlarged by the Normans, used by the Hohenstaufen
     and then the Angevins."@en ;
  geo:lat "38.1115"^^xsd:decimal ;
  geo:long "13.3523"^^xsd:decimal ;
  ohm:period v1282:period_1282 ;
  ohm:buildingType "palazzo"@it , "palace"@en ;
  v1282:narrativeRole "Sede di Jean de Saint-Rémy"@it ;
  owl:sameAs <https://www.wikidata.org/entity/Q848055> .

ohmplace:palermo_cattedrale a ohm:HistoricalBuilding ;
  rdfs:label "Cattedrale di Palermo"@it ,
             "Cathedral of Palermo"@en ;
  dcterms:description
    "Sede dell'arcivescovo di Palermo. Fondata dai Normanni
     su una precedente moschea araba."@it ,
    "Seat of the Archbishop of Palermo. Founded by the Normans
     on a previous Arab mosque."@en ;
  geo:lat "38.1129"^^xsd:decimal ;
  geo:long "13.3517"^^xsd:decimal ;
  ohm:period v1282:period_1282 ;
  ohm:buildingType "cattedrale"@it , "cathedral"@en ;
  v1282:narrativeRole "Sede dell'arcivescovo — fazione Clero"@it ;
  owl:sameAs <https://www.wikidata.org/entity/Q847497> .

ohmplace:palermo_porto a ohm:HistoricalPlace ;
  rdfs:label "Porto di Palermo"@it , "Port of Palermo"@en ;
  dcterms:description
    "Porto principale di Palermo. Nel 1282 sede della flotta
     di Carlo d'Angiò in allestimento per la spedizione
     contro Bisanzio."@it ,
    "Main port of Palermo. In 1282, site of Charles of Anjou's
     fleet being assembled for the expedition against Byzantium."@en ;
  geo:lat "38.1200"^^xsd:decimal ;
  geo:long "13.3700"^^xsd:decimal ;
  ohm:period v1282:period_1282 ;
  v1282:narrativeRole "Flotta angioina — punto critico"@it .

ohmplace:palermo_mercato a ohm:HistoricalPlace ;
  rdfs:label "Mercato della Kalsa"@it , "Kalsa Market"@en ;
  dcterms:description
    "Mercato principale del quartiere Kalsa. Centro della
     vita quotidiana e delle voci della città."@it ,
    "Main market of the Kalsa quarter. Centre of daily life
     and city rumours."@en ;
  geo:lat "38.1138"^^xsd:decimal ;
  geo:long "13.3660"^^xsd:decimal ;
  ohm:period v1282:period_1282 ;
  v1282:narrativeRole "Informazioni — fazione Popolo"@it .
\end{verbatim}

\section{Altri Centri Siciliani}

\begin{verbatim}
ohmplace:messina_1282 a ohm:HistoricalCity ;
  rdfs:label "Messina (1282)"@it , "Messina (1282)"@en ;
  dcterms:description
    "Secondo centro della Sicilia. La rivolta si estese
     a Messina dopo i Vespri di Palermo. Porta verso
     la Calabria e il continente."@it ,
    "Second city of Sicily. The uprising spread to Messina
     after the Palermo Vespers. Gateway to Calabria and
     the mainland."@en ;
  geo:lat "38.1938"^^xsd:decimal ;
  geo:long "15.5540"^^xsd:decimal ;
  ohm:period v1282:period_1282 ;
  v1282:narrativeRole "Canale verso la Sicilia orientale"@it ;
  owl:sameAs <https://www.wikidata.org/entity/Q16"> .

ohmplace:corleone_1282 a ohm:HistoricalCity ;
  rdfs:label "Corleone (1282)"@it , "Corleone (1282)"@en ;
  dcterms:description
    "Centro dell'entroterra siciliano. Partecipò attivamente
     alla rivolta dopo il 30 marzo."@it ,
    "Town in the Sicilian interior. Actively participated in
     the uprising after 30 March."@en ;
  geo:lat "37.8145"^^xsd:decimal ;
  geo:long "13.0693"^^xsd:decimal ;
  ohm:period v1282:period_1282 .

ohmplace:trapani_1282 a ohm:HistoricalCity ;
  rdfs:label "Trapani (1282)"@it , "Trapani (1282)"@en ;
  dcterms:description
    "Porto occidentale della Sicilia. Punto di sbarco
     potenziale per truppe aragonesi."@it ,
    "Western port of Sicily. Potential landing point
     for Aragonese forces."@en ;
  geo:lat "38.0176"^^xsd:decimal ;
  geo:long "12.5140"^^xsd:decimal ;
  ohm:period v1282:period_1282 .

ohmplace:siracusa_1282 a ohm:HistoricalCity ;
  rdfs:label "Siracusa (1282)"@it , "Syracuse (1282)"@en ;
  dcterms:description
    "Centro storico della Sicilia sud-orientale.
     Antica capitale greca e bizantina."@it ,
    "Historic centre of south-eastern Sicily.
     Ancient Greek and Byzantine capital."@en ;
  geo:lat "37.0755"^^xsd:decimal ;
  geo:long "15.2866"^^xsd:decimal ;
  ohm:period v1282:period_1282 .

ohmplace:agrigento_1282 a ohm:HistoricalCity ;
  rdfs:label "Agrigento (1282)"@it , "Agrigento (1282)"@en ;
  dcterms:description
    "Centro del versante meridionale. Sede vescovile."@it ,
    "Centre of the southern coast. Episcopal seat."@en ;
  geo:lat "37.3111"^^xsd:decimal ;
  geo:long "13.5765"^^xsd:decimal ;
  ohm:period v1282:period_1282 .

ohmplace:catania_1282 a ohm:HistoricalCity ;
  rdfs:label "Catania (1282)"@it , "Catania (1282)"@en ;
  dcterms:description
    "Secondo porto della Sicilia orientale, ai piedi
     dell'Etna."@it ,
    "Second port of eastern Sicily, at the foot of Etna."@en ;
  geo:lat "37.5023"^^xsd:decimal ;
  geo:long "15.0873"^^xsd:decimal ;
  ohm:period v1282:period_1282 .
\end{verbatim}

\section{Itinerari Narrativi}

\begin{verbatim}
%% Itinerario: La processione di Santo Spirito

v1282:route_processione a ohm:HistoricalRoute ;
  rdfs:label "Processione di Santo Spirito — 30 marzo 1282"@it ,
             "Santo Spirito Procession — 30 March 1282"@en ;
  ohm:startPoint ohmplace:palermo_cattedrale ;
  ohm:endPoint ohmplace:palermo_santo_spirito ;
  ohm:period v1282:period_1282 ;
  dcterms:description
    "Percorso della processione pasquale dalla cattedrale
     alla chiesa di Santo Spirito, teatro dei Vespri."@it .

%% Itinerario: Rete dei congiurati

v1282:network_congiurati a ohm:HistoricalNetwork ;
  rdfs:label "Rete dei Congiurati"@it ,
             "Conspirators' Network"@en ;
  ohm:node ohmplace:palermo_kalsa ,
           ohmplace:messina_1282 ,
           ohmplace:trapani_1282 ;
  ohm:period v1282:period_1282 ;
  dcterms:description
    "Rete di comunicazione clandestina della cospirazione
     aragonese in Sicilia."@it .
\end{verbatim}

\section{Note Metodologiche}

Le coordinate geografiche di questa appendice sono approssimazioni
storico-topografiche. Palermo medievale aveva una planimetria
diversa da quella odierna, ma il nucleo del Cassaro e la posizione
della Kalsa rispetto al porto sono sostanzialmente confermati
dalla storiografia e dall'archeologia urbana.

Per integrare questi dati in OpenHistoryMap, importare il file
Turtle in un endpoint SPARQL compatibile con GeoSPARQL.
I dati sono rilasciati sotto licenza \textbf{CC BY 4.0}.

\filettodoppio
